%============================================================
\section{Discussion}
%============================================================
\paragraph{What matters for IRENE-style fusion at small scale.}
Three findings are consistent across metrics and seeds.
(1)~\emph{Fusion helps, and the text representation decides by how much.} Adding the note
and structured data to the image raises macro AUROC by about 0.09--0.10. However, a weak note
encoder (static word vectors alone) wiped out most of this gain. A document-level token
recovered it.
(2)~\emph{Keep intra-modal self-attention, and limit cross-attention to early layers.}
IRENE's (self+cross)/2 blocks followed by joint self-attention are on par with early and
late fusion. Cross-attention without self-attention, or cross-attention at every layer, is
clearly worse. On a few hundred patients, extra cross-modal capacity causes overfitting.
(3)~\emph{Contrastive alignment is not free.} As an auxiliary loss on hundreds of pairs,
CLIP and SigLIP hurt a mid-fusion model and are neutral for a dual encoder. The benefit reported
for medical vision--language models comes from large-scale pre-training, which our setting
lacks. A natural next step is to initialise the two streams from a publicly released
radiograph--report encoder (e.g.\ BiomedCLIP or CXR-CLIP) instead of learning the alignment
on the target data.

\paragraph{What the headline numbers hide.} The dataset is public, multimodal and, with IRENE's
input types, a natural testbed. It is also strongly confounded by source. A source-only probe
matches the fusion transformers, and a linear bag-of-words model on the masked notes beats all
of them. The de-confounded subset preserves the ranking of the fusion strategies, which supports
the architectural conclusions. The absolute numbers, however, should not be read as clinical
diagnostic accuracy.

%============================================================
\section{Limitations}
\label{sec:limits}
%============================================================
\begin{itemize}
\item \textbf{Scale.} \NImages{} images from \NPatients{} patients, and three or fewer positives
per test fold for some classes. We pool out-of-fold predictions and report bootstrap intervals.
Many ablation differences are nevertheless within noise, and the ablations use 2 seeds.
\item \textbf{Compute.} No GPU was available. The model is a narrow IRENE ($d{=}128$, 6 layers,
16 image tokens from a frozen CNN) rather than the original ViT-B/16 initialised from ImageNet
weights, and the text encoder is not a clinical BERT. The recipe was revised once, after a full
cross-validation pilot showed that the note branch was the bottleneck (it added the document token
and class-balanced loss). This single round of development used out-of-fold results, so a fresh
external test set would be needed for an unbiased estimate.
\item \textbf{Data.} The labels come from publication metadata, the notes are written by case
authors rather than being contemporaneous EHR text, and the source confounding described in
\cref{sec:leak} limits clinical interpretation. The engineering code supports large sharded
image--text corpora and multi-GPU training, but this study did not exercise that scale.
\end{itemize}

%============================================================
\section{Conclusion}
%============================================================
We provide an open, end-to-end re-implementation of the IRENE multimodal transformer on
public data. It includes a documented image--text--structured pipeline, a modern PyTorch
training stack, a single model that covers early, mid and late fusion, six attention structures
and two contrastive objectives, and a controlled study with patient-grouped cross-validation.
The results support IRENE's core design: shallow bidirectional attention that retains
self-attention, followed by joint layers. They also show that a competent text representation,
robustness to missing notes and awareness of dataset confounding matter more than the
particular choice of fusion operator.
